\section{Analytic framework}

The argument has two analytic joints.  Products and Fourier multipliers must
keep every differentiated velocity--pressure row in a common Sobolev scale;
at the terminal time, an \(H^3\) local theory must restart that same
pressure-normalized history.  We record the estimates needed at those two
joints, with the constants later used in the quantitative bounds.

\begin{definition}
For \(f\in\mathcal S(\R^3)\), set
\[
 \widehat f(\xi):=(2\pi)^{-3/2}\int_{\R^3}e^{-ix\cdot\xi}f(x)\dd x,
 \qquad
 \langle\xi\rangle:=(1+\abs{\xi}^2)^{1/2}.
\]
This integral defines \(\widehat f\) for \(f\in L^1(\R^3)\).  The
Fourier transform on \(\mathcal S'(\R^3)\) is its extension by duality.
For \(s\in\R\), define
\[
 H^s(\R^3):=
 \left\{f\in\mathcal S'(\R^3):
       \langle\xi\rangle^s\widehat f\in L^2(\R^3)\right\},
\]
and define the finite homogeneous-norm domain by
\[
 \dot H^s(\R^3):=
 \left\{f\in\mathcal S'(\R^3):
 \begin{array}{l}
 \widehat f\text{ is represented by a measurable function on }\R^3,\\[-2pt]
 \abs\xi^s\widehat f\in L^2(\R^3)
 \end{array}
 \right\}.
\]
On these domains set
\[
 \norm f_{H^s}^2:=\int_{\R^3}\langle\xi\rangle^{2s}
 \abs{\widehat f(\xi)}^2\dd\xi,
 \qquad
 \norm f_{\dot H^s}^2:=\int_{\R^3}\abs{\xi}^{2s}
 \abs{\widehat f(\xi)}^2\dd\xi.
\]
Set \(A:=(1-\Delta)^{1/2}\), \(\Lambda:=(-\Delta)^{1/2}\),
\[
 H^s_\sigma:=\{v\in H^s(\R^3;\R^3):\nabla\cdot v=0\},
 \qquad
 H^\infty:=\bigcap_{m\in\Nzero}H^m,
 \qquad
 H^\infty_\sigma:=\bigcap_{m\in\Nzero}H^m_\sigma.
\]
\[
 \mathcal S_\sigma(\R^3):=
 \{v\in\mathcal S(\R^3;\R^3):\nabla\cdot v=0\}.
\]
For a vector or tensor \(F=(F_I)_I\), all Sobolev norms are componentwise:
\[
 \norm F_{H^s}^2:=\sum_I\norm{F_I}_{H^s}^2,
 \qquad
 \norm F_{\dot H^s}^2:=\sum_I\norm{F_I}_{\dot H^s}^2.
\]
Repeated Latin indices are summed from \(1\) to \(3\).
\end{definition}

\begin{lemma}[Sobolev nesting and embedding]
\label{lem:sobolev-embedding}
If \(s_2\ge s_1\), then
\[
 \norm f_{H^{s_1}}\le\norm f_{H^{s_2}}.
\]
If \(k\in\Nzero\) and \(s>k+3/2\), then
\[
 \max_{\abs\alpha\le k}\norm{\partial^\alpha f}_{L^\infty}
 \le C_{s,k}^{\rm emb}\norm f_{H^s},
\]
where
\[
 C_{s,k}^{\rm emb}:=(2\pi)^{-3/2}
 \max_{0\le j\le k}
 \left(\int_{\R^3}\abs\xi^{2j}\langle\xi\rangle^{-2s}\dd\xi\right)^{1/2}<\infty.
\]
In particular,
\[
 H^\infty(\R^3)\hookrightarrow C_b^\infty(\R^3).
\]
\end{lemma}

\begin{proof}
The first assertion follows from
\(\langle\xi\rangle^{2s_1}\le\langle\xi\rangle^{2s_2}\).
For \(\abs\alpha\le k\), Fourier inversion and Cauchy--Schwarz give
\[
 \abs{\partial^\alpha f(x)}
 \le(2\pi)^{-3/2}
 \left(\int_{\R^3}\abs\xi^{2\abs\alpha}
 \langle\xi\rangle^{-2s}\dd\xi\right)^{1/2}
 \norm f_{H^s}.
\]
The integral is finite exactly when \(s>\abs\alpha+3/2\).
Approximation by Schwartz functions and dominated convergence give continuous
representatives.  For each \(k\), choose \(s=k+2\).
\end{proof}

\begin{lemma}[Explicit three-dimensional homogeneous Sobolev constants]
\label{lem:explicit-homogeneous-sobolev}
For \(3\le p\le6\), set
\[
 s(p):=3\left(\frac12-\frac1p\right)
\]
and
\begin{equation}
 \mathsf S(p):=
 2^{-s(p)}\pi^{-s(p)/2}
 \left(
  \frac{\Gamma((3-2s(p))/2)}
       {\Gamma((3+2s(p))/2)}
 \right)^{1/2}
 \left(\frac{\Gamma(3)}{\Gamma(3/2)}\right)^{s(p)/3}.
 \label{eq:explicit-homogeneous-sobolev-constant}
\end{equation}
Then every scalar, vector, or finite tensor field \(F\) in
\(\dot H^{s(p)}(\R^3)\) satisfies
\begin{equation}
 \norm F_{L^p(\R^3)}
 \le \mathsf S(p)\norm F_{\dot H^{s(p)}(\R^3)}.
 \label{eq:explicit-homogeneous-sobolev}
\end{equation}
In particular,
\[
 \mathsf S(3)=2^{-1/6}\pi^{-1/3},
 \qquad
 \mathsf S(6)=\frac{2^{2/3}}{\sqrt3\,\pi^{2/3}},
 \qquad
 \mathsf S(6)\mathsf S(3)=\frac1\pi\sqrt{\frac23}.
\]
\end{lemma}

\begin{proof}
For scalar fields, \eqref{eq:explicit-homogeneous-sobolev} is the sharp
three-dimensional homogeneous Sobolev inequality obtained from the sharp
Hardy--Littlewood--Sobolev inequality \cite{Lieb1983}.  After duality,
insertion of the Riesz-potential kernel constant, and conversion to the
unitary Fourier transform fixed above, the squared sharp Sobolev constant for
the multiplier
\(\widehat{\Lambda^sf}(\xi)=|\xi|^s\widehat f(\xi)\) is
\begin{equation}
 2^{-2s}\pi^{-s}
 \frac{\Gamma((3-2s)/2)}{\Gamma((3+2s)/2)}
 \left(\frac{\Gamma(3)}{\Gamma(3/2)}\right)^{2s/3}.
 \label{eq:explicit-homogeneous-sobolev-squared-constant}
\end{equation}
Taking the square root in
\eqref{eq:explicit-homogeneous-sobolev-squared-constant} and putting
\(p=6/(3-2s)\) gives
\eqref{eq:explicit-homogeneous-sobolev-constant}--
\eqref{eq:explicit-homogeneous-sobolev}.  For a finite component array
\(F=(F_I)_I\), Minkowski's inequality in \(L^{p/2}\) gives
\[
 \norm F_p^2
 =\left\|\sum_I|F_I|^2\right\|_{p/2}
 \le\sum_I\norm{F_I}_p^2
 \le\mathsf S(p)^2\sum_I\norm{F_I}_{\dot H^{s(p)}}^2.
\]
The two endpoint formulas follow by evaluating the gamma factors at
\(s=1/2\) and \(s=1\).
\end{proof}

\Needspace{34\baselineskip}
\begin{lemma}[Sobolev multiplication]
\label{lem:sobolev-product}
Set
\[
 \mathsf H_4:=(2\pi)^{-3/4}
 \left(\int_{\R^3}\langle\xi\rangle^{-4}\dd\xi\right)^{1/4},
 \qquad C_0^{\rm alg}:=\mathsf H_4^2.
\]
For \(m\ge1\), define
\begin{align*}
 \mathsf E_m
 &:=\max_{0\le j\le\lfloor m/2\rfloor}C_{j+2,j}^{\rm emb},
 &\mathsf C_m^{\rm mix}&:=13^{m/2}\mathsf E_m,\\
 \mathsf C_m^{\rm F}
 &:=2^m(2\pi)^{-3/2}
 \left(\int_{\R^3}\langle\xi\rangle^{-2m}\dd\xi\right)^{1/2}
 &&(m\ge2),
\end{align*}
and set
\[
 C_1^{\rm alg}:=\mathsf C_1^{\rm mix},
 \qquad
 C_m^{\rm alg}:=\max\{\mathsf C_m^{\rm mix},\mathsf C_m^{\rm F}\}
 \quad(m\ge2).
\]
All these constants are finite and explicit.  For every \(m\in\Nzero\),
\[
 \norm{fg}_{H^m}
 \le C_m^{\rm alg}\norm f_{H^{m+1}}\norm g_{H^{m+1}}.
\]
For \(m\ge2\),
\[
 \norm{fg}_{H^m}
 \le C_m^{\rm alg}\norm f_{H^m}\norm g_{H^m}.
\]
\end{lemma}

\begin{proof}
Hausdorff--Young and H\"older give
\[
 \norm f_4
 \le(2\pi)^{-3/4}\norm{\widehat f}_{4/3}
 \le\mathsf H_4\norm f_{H^1}.
\]
Thus \(\norm{fg}_2\le C_0^{\rm alg}\norm f_{H^1}\norm g_{H^1}\).

Let \(m\ge1\).  The exact integer-order identity
\[
 \norm h_{H^m}^2
 =\sum_{\abs\alpha\le m}
 \frac{m!}{(m-\abs\alpha)!\,\alpha!}
 \norm{\partial^\alpha h}_2^2
\]
follows by expanding \((1+\abs\xi^2)^m\).  In each Leibniz term for
\(\partial^\alpha(fg)\), place the factor carrying at most
\(\lfloor\abs\alpha/2\rfloor\) derivatives in \(L^\infty\).  The definition
of \(\mathsf E_m\), Sobolev nesting, and
\(\sum_{\beta\le\alpha}\binom\alpha\beta=2^{\abs\alpha}\) give
\[
 \norm{\partial^\alpha(fg)}_2
 \le \mathsf E_m2^{\abs\alpha}
 \norm f_{H^{m+1}}\norm g_{H^{m+1}}.
\]
The multinomial identity
\[
 \sum_{\abs\alpha\le m}
 \frac{m!}{(m-\abs\alpha)!\,\alpha!}4^{\abs\alpha}=13^m
\]
proves the first estimate with \(\mathsf C_m^{\rm mix}\).

For \(m\ge2\), the inequalities
\[
 \langle\eta+\zeta\rangle^m
 \le2^{m-1}\bigl(\langle\eta\rangle^m+
                         \langle\zeta\rangle^m\bigr),
 \qquad
 \norm{\widehat f}_1
 \le\left(\int_{\R^3}\langle\xi\rangle^{-2m}\dd\xi\right)^{1/2}
       \norm f_{H^m},
\]
combined with \(\widehat{fg}=(2\pi)^{-3/2}\widehat f*\widehat g\)
and Young's convolution inequality give
\[
 \norm{fg}_{H^m}
 \le\mathsf C_m^{\rm F}\norm f_{H^m}\norm g_{H^m}.
\]
The chosen maximum proves both asserted estimates with this constant.
\end{proof}

\begin{lemma}[Riesz and Leray multipliers]
\label{lem:riesz-leray}
Let
\[
 r_j(\xi):=
 \begin{cases}
 i\xi_j/\abs\xi,&\xi\ne0,\\
 0,&\xi=0,
 \end{cases}
 \qquad
 P(\xi):=
 \begin{cases}
 I-\dfrac{\xi\otimes\xi}{\abs\xi^2},&\xi\ne0,\\
 I,&\xi=0,
 \end{cases}
\]
and
\[
 \widehat{R_jf}(\xi)=r_j(\xi)\widehat f(\xi),
 \qquad
 \widehat{\Pp f}(\xi)=P(\xi)\widehat f(\xi).
\]
For all \(s\in\R\),
\[
 \norm{R_jf}_{H^s}\le\norm f_{H^s},
 \qquad
 \norm{\Pp f}_{H^s}\le\norm f_{H^s}.
\]
If \(m\in\Nzero\) and \(v\in H^{m+1}_\sigma\), then
\[
 q:=\sum_{i,j=1}^3R_iR_j(v_iv_j)\in H^m,
 \qquad
 -\Delta q=\sum_{i,j=1}^3\partial_i\partial_j(v_iv_j),
\]
and
\[
 \Pp\nabla\cdot(v\otimes v)
 =(v\cdot\nabla)v+\nabla q.
\]
Here
\[
 (v\otimes v)_{ij}:=v_iv_j,
 \qquad
 (\nabla\cdot(v\otimes v))_i
 :=\sum_{j=1}^3\partial_j(v_iv_j),
 \qquad
 ((v\cdot\nabla)v)_i:=\sum_{j=1}^3v_j\partial_jv_i.
\]
\end{lemma}

\begin{proof}
The multiplier symbols have operator norm at most one.  Plancherel proves the
first two estimates.  Lemma~\ref{lem:sobolev-product} gives
\(v_iv_j\in H^m\).  The remaining identities follow after multiplication by
\(\abs\xi^2\) and by the symbol of \(\Pp\).
\end{proof}

\begin{lemma}[Mixed Leibniz identities]
\label{lem:mixed-leibniz}
For \(N\in\Nzero\) and \(\alpha\in\Nzero^3\),
\[
 \partial_t^N\partial_x^\alpha(fg)
 =\sum_{a=0}^N\sum_{\beta\le\alpha}
 \binom Na\binom\alpha\beta
 (\partial_t^a\partial_x^\beta f)
 (\partial_t^{N-a}\partial_x^{\alpha-\beta}g).
\]
\end{lemma}

\begin{proof}
Apply the one-variable Leibniz identity successively in \(t,x_1,x_2,x_3\).
\end{proof}

\Needspace{14\baselineskip}
\begin{lemma}[Adjacent Hilbert-scale trace]
\label{lem:adjacent-trace}
Let \(0<T<\infty\), \(m\in\R\), and
\[
 v\in L^2(0,T;H^{m+1}),
 \qquad
 \partial_tv\in L^2(0,T;H^{m-1}).
\]
Then \(v\) has a unique representative in \(C([0,T];H^m)\), and, for
\(0\le r\le t\le T\),
\begin{align}
 \norm{v(t)}_{H^m}^2-\norm{v(r)}_{H^m}^2
 &=2\operatorname{Re}\int_r^t
 \ip{A^{m-1}\partial_tv(s)}{A^{m+1}v(s)}_{L^2}\dd s,
 \label{eq:adjacent-trace-identity}\\
 \abs{\norm{v(t)}_{H^m}^2-\norm{v(r)}_{H^m}^2}
 &\le2\norm{\partial_tv}_{L^2(r,t;H^{m-1})}
 \norm v_{L^2(r,t;H^{m+1})},
 \label{eq:adjacent-trace-modulus}\\
 \sup_{0\le t\le T}\norm{v(t)}_{H^m}^2
 &\le T^{-1}\norm v_{L^2(0,T;H^m)}^2
 +2\norm{\partial_tv}_{L^2(0,T;H^{m-1})}
 \norm v_{L^2(0,T;H^{m+1})}.
 \label{eq:adjacent-trace-sup}
\end{align}
\end{lemma}

\begin{proof}
Let \(S_R\) denote Fourier restriction to \(\{\abs\xi\le R\}\).  For
\(v_R=S_Rv\),
\[
 \frac{\dd}{\dd t}\norm{v_R(t)}_{H^m}^2
 =2\operatorname{Re}\ip{A^{m-1}\partial_tv_R(t)}
 {A^{m+1}v_R(t)}_{L^2}.
\]
Integration and Cauchy--Schwarz prove
\eqref{eq:adjacent-trace-identity}--\eqref{eq:adjacent-trace-modulus} for
\(v_R\).  Applied to \(S_Rv-S_Sv\), this estimate shows that
\((v_R)_R\) is Cauchy in \(C([0,T];H^m)\); choose one Lebesgue time for the
initial Cauchy term.  Its limit equals \(v\) in
\(L^2(0,T;H^m)\).  Passage to the limit proves the first two formulas.
Choose \(r\in(0,T)\) with
\(\norm{v(r)}_{H^m}^2\le T^{-1}\norm v_{L^2H^m}^2\) and apply
\eqref{eq:adjacent-trace-modulus}.
\end{proof}

\begin{lemma}[Explicit integer transport commutator]
\label{lem:transport-commutator}
Set
\begin{align}
 \Gamma_{1,1}&:=1,\notag\\
 \Gamma_{\ell,k}
 &:=(2\ell-4+\sqrt3)^{(k-1)(\ell-k)}
 \qquad(\ell\ge2,\ 1\le k\le\ell),
 \label{eq:explicit-Gamma-lk}\\
 K_m&:=
 \sum_{\ell=1}^{m}3^\ell\binom m\ell
 \sum_{k=1}^{\ell}\binom\ell k\,\Gamma_{\ell,k}
 \qquad(m\ge3).
 \label{eq:explicit-commutator-constant}
\end{align}
Every \(v\in H^{m+1}_\sigma(\R^3)\) satisfies
\begin{equation}
 \abs{\operatorname{Re}
 \ip{A^m\Pp\nabla\cdot(v\otimes v)}{A^mv}_{L^2}}
 \le K_m\norm{\nabla v}_{L^\infty}\norm v_{H^m}^2.
 \label{eq:transport-commutator}
\end{equation}
In particular,
\begin{equation}
 K_m\le(7^m-1)(2m-4+\sqrt3)^{(m-1)^2/4}.
 \label{eq:explicit-commutator-upper-bound}
\end{equation}
\end{lemma}

\begin{proof}
We first prove the interpolation inequality used in the commutator.  For the
full ordered derivative tensor \(\nabla^jF\), fix \(M\ge1\) and set
\[
 p_{M,0}:=\infty,
 \qquad
 p_{M,j}:=\frac{2M}{j}\quad(1\le j\le M),
 \qquad
 A_M:=2M-2+\sqrt3.
\]
Then
\begin{equation}
 \norm{\nabla^jF}_{L^{p_{M,j}}}
 \le A_M^{j(M-j)/2}
 \norm F_\infty^{\,1-j/M}
 \norm{\nabla^MF}_2^{\,j/M}
 \qquad(0\le j\le M).
 \label{eq:explicit-derivative-interpolation}
\end{equation}
The endpoint cases are identities.  For \(1\le j\le M-1\), put
\[
 G:=\nabla^{j-1}F,
 \qquad W:=\nabla G,
 \qquad p:=p_{M,j}.
\]
All contractions below are the Euclidean contractions of the corresponding
ordered tensors.  Integration by parts gives
\begin{align*}
 \int_{\R^3}|W|^p
 ={}&-\int_{\R^3}|W|^{p-2}G:\Delta G\\
 &-(p-2)\int_{\R^3}|W|^{p-4}
   \sum_{q=1}^3(G:W_{\cdot q})(W:\partial_qW).
\end{align*}
The pointwise bounds
\[
 |\Delta G|\le\sqrt3\,|\nabla^2G|,
 \qquad
 \left|\sum_{q=1}^3(G:W_{\cdot q})(W:\partial_qW)\right|
 \le |G|\,|W|^2|\nabla^2G|
\]
give
\[
 \int_{\R^3}|W|^p
 \le(p-2+\sqrt3)
 \int_{\R^3}|G|\,|W|^{p-2}|\nabla^2G|.
\]
Since
\[
 \frac1{p_{M,j-1}}+\frac{p-2}{p}
 +\frac1{p_{M,j+1}}=1,
\]
H\"older's inequality, with
\(
 X_j:=\norm{\nabla^jF}_{L^{p_{M,j}}}
\), gives
\begin{equation}
 X_j^2
 \le\left(\frac{2M}{j}-2+\sqrt3\right)X_{j-1}X_{j+1}
 \le A_MX_{j-1}X_{j+1}.
 \label{eq:explicit-interpolation-three-point}
\end{equation}
The numbers \(d_j:=j(M-j)/2\) obey
\[
 2d_j-d_{j-1}-d_{j+1}=1,
 \qquad d_0=d_M=0.
\]
The discrete maximum principle applied to
\eqref{eq:explicit-interpolation-three-point} gives
\[
 X_j\le A_M^{d_j}X_0^{1-j/M}X_M^{j/M},
\]
which is \eqref{eq:explicit-derivative-interpolation}.  The integration by
parts is justified first with
\((|W|^2+\varepsilon)^{(p-2)/2}\), followed by
\(\varepsilon\downarrow0\); smooth approximation handles zero tensor norms.

For \(\abs\alpha\le m\), set
\[
 \alpha!:=\alpha_1!\alpha_2!\alpha_3!,
 \qquad
 c_{m,\alpha}:=\frac{m!}{(m-\abs\alpha)!\,\alpha!}.
\]
The multiplier \(A^{2m}=(1-\Delta)^m\) and the multinomial identity give the
exact pairing formula
\begin{equation}
 \ip{A^mf}{A^mg}_{L^2}
 =\sum_{\abs\alpha\le m}c_{m,\alpha}
   \ip{\partial^\alpha f}{\partial^\alpha g}_{L^2}.
 \label{eq:integer-Am-pairing}
\end{equation}
The Leray projection is self-adjoint, commutes with \(A^m\), and fixes
\(A^mv\).  Since
\(\nabla\cdot(v\otimes v)=(v\cdot\nabla)v\), it follows that
\begin{equation}
 \ip{A^m\Pp\nabla\cdot(v\otimes v)}{A^mv}_{L^2}
 =\sum_{\abs\alpha\le m}c_{m,\alpha}
   \ip{\partial^\alpha((v\cdot\nabla)v)}
      {\partial^\alpha v}_{L^2}.
 \label{eq:integer-projected-pairing}
\end{equation}
For each \(\alpha\),
\[
 \partial^\alpha((v\cdot\nabla)v)
 =(v\cdot\nabla)\partial^\alpha v
 +\sum_{0<\beta\le\alpha}\binom\alpha\beta
 (\partial^\beta v\cdot\nabla)
 \partial^{\alpha-\beta}v.
\]
The divergence-free transport term has zero real pairing with
\(\partial^\alpha v\), because
\[
 \operatorname{Re}\ip{(v\cdot\nabla)\partial^\alpha v}
 {\partial^\alpha v}_{L^2}
 =\frac12\int_{\R^3}v\cdot\nabla
   \abs{\partial^\alpha v}^2\dd x=0.
\]
\begin{samepage}
For a remaining term, put
\[
 \ell:=\abs\alpha,
 \qquad k:=\abs\beta.
\]
\par
\end{samepage}
\begin{samepage}
If \(\ell\ge2\), apply
\eqref{eq:explicit-derivative-interpolation} to \(F=\nabla v\), with
\[
 M=\ell-1,
 \qquad j=k-1.
\]
The two H\"older exponents are \(p_{M,j}\) and \(p_{M,M-j}\), whose
reciprocal sum is \(1/2\).  The two factors are subarrays of the corresponding
ordered derivative tensors, so
\begin{align}
 \norm{(\partial^\beta v\cdot\nabla)
          \partial^{\alpha-\beta}v}_2
 &\le
 \norm{\partial^\beta v}_{L^{p_{M,j}}}
 \norm{\nabla\partial^{\alpha-\beta}v}_{L^{p_{M,M-j}}}\notag\\
 &\le\Gamma_{\ell,k}\norm{\nabla v}_\infty
       \norm{\nabla^\ell v}_2.
 \label{eq:explicit-commutator-product}
\end{align}
\par
\end{samepage}
\begin{samepage}
For \(\ell=1\), this formula follows from the direct
\(L^\infty\)-\(L^2\) estimate and \(\Gamma_{1,1}=1\).  Since
\[
 \norm{\nabla^\ell v}_2=\norm v_{\dot H^\ell}\le\norm v_{H^m},
 \qquad
 \norm{\partial^\alpha v}_2\le\norm v_{H^m},
\]
each remaining pairing is bounded by
\[
 \Gamma_{\ell,k}\norm{\nabla v}_\infty\norm v_{H^m}^2.
\]
\par
\end{samepage}
The multinomial identities
\[
 \sum_{\abs\alpha=\ell}c_{m,\alpha}
 =3^\ell\binom m\ell,
 \qquad
 \sum_{\substack{\beta\le\alpha\\\abs\beta=k}}
 \binom\alpha\beta=\binom\ell k
\]
now give \eqref{eq:transport-commutator} with the constant
\eqref{eq:explicit-commutator-constant}.  Finally,
\[
 \sum_{k=1}^{\ell}\binom\ell k\le2^\ell,
 \qquad
 (k-1)(\ell-k)\le\frac{(\ell-1)^2}{4},
\]
and summation of
\(
 \binom m\ell3^\ell2^\ell
\)
over \(1\le\ell\le m\) proves
\eqref{eq:explicit-commutator-upper-bound}.  Approximation by divergence-free
Schwartz fields justifies all integrations by parts for
\(v\in H^{m+1}_\sigma\).
\end{proof}

The preceding estimates control the differentiated rows as they approach an
endpoint.  To continue their common limit as the original solution, we now fix a
time-translation-invariant local class that admits both late preterminal
slices and the endpoint restart, and in which uniqueness holds on every
overlap.

\begin{definition}[Shifted pressure-normalized \(H^3\) mild class]
\label{def:shifted-normalized-H3-mild-class}
Let \(0\le s<b<\infty\) and \(a\in H^3_\sigma(\R^3)\).  Define
\(\mathscr M_\nu^3(s,b;a)\) to be the set of pairs \((w,q)\) such that
\begin{align}
 w&\in C([s,b];H^3_\sigma(\R^3)),
 &q&\in C([s,b];H^3(\R^3)),
 \label{eq:shifted-mild-class-spaces}\\
 w(s)&=a,
 &q(t)&=R_iR_j\bigl(w_i(t)w_j(t)\bigr),
 \label{eq:shifted-mild-class-datum-pressure}
\end{align}
The remaining defining condition is the Duhamel identity: for every
\(t\in[s,b]\),
\begin{equation}
 w(t)=e^{\nu(t-s)\Delta}a
 -\int_s^t e^{\nu(t-r)\Delta}
 \Pp\nabla\!\cdot(w\otimes w)(r)\dd r
 \quad\text{in }H^3_\sigma(\R^3).
 \label{eq:shifted-mild-class-duhamel}
\end{equation}
The integral is an \(H^3_\sigma\)-valued Bochner integral: multiplication in
\(H^3(\R^3)\) places
\(\Pp\nabla\!\cdot(w\otimes w)\) in \(C([s,b];H^2_\sigma)\), and the heat
multiplier from \(H^2\) to \(H^3\) has an integrable
\((t-r)^{-1/2}\) singularity.
The maximal lifespan in this class is
\begin{equation}
 \tau_{\nu,3}^{\max}(s,a)
 :=\sup\left\{b-s:\ b>s,\
 \mathscr M_\nu^3(s,b;a)\ne\varnothing\right\}\in[0,\infty].
 \label{eq:shifted-mild-maximal-lifespan}
\end{equation}
\end{definition}

\begin{lemma}[Restriction and classical admission]
\label{lem:shifted-mild-restriction}
If \((w,q)\in\mathscr M_\nu^3(s,b;a)\) and
\(s\le r<c\le b\), then
\begin{equation}
 (w,q)|_{[r,c]}\in\mathscr M_\nu^3(r,c;w(r)).
 \label{eq:shifted-mild-restriction}
\end{equation}
If \((u,p)\) is a pressure-complete classical history in the sense of
Definition~\ref{def:formal-maximal-history} and
\([s,b]\) is contained in its interval of existence, then
\begin{equation}
 (u,p)|_{[s,b]}\in\mathscr M_\nu^3(s,b;u(s)).
 \label{eq:classical-mild-admission}
\end{equation}
\end{lemma}

\begin{proof}
Split \eqref{eq:shifted-mild-class-duhamel} at \(r\) and use the heat
semigroup identity.  This gives, for \(r\le t\le c\),
\[
 w(t)=e^{\nu(t-r)\Delta}w(r)
 -\int_r^t e^{\nu(t-\rho)\Delta}
 \Pp\nabla\!\cdot(w\otimes w)(\rho)\dd\rho.
\]
The pressure formula is unchanged under restriction.  For a classical
history, apply \(\Pp\) to the momentum equation and use variation of
constants; the decaying pressure formula gives
\eqref{eq:shifted-mild-class-datum-pressure}.
\end{proof}

\begin{theorem}[Quantitative local restart]
\label{thm:local-restart}
Set \(C_{\rm q}:=\sqrt3\,C_3^{\rm alg}\).  For \(\nu>0\), set
\[
 L_\nu:=4(1+\nu+\nu^{-1})
 \exp\!\left(\frac{\nu+\nu^{-1}}2\right),
 \qquad
 c_\nu:=(8C_{\rm q}L_\nu^2)^{-2}.
\]
For every \(s\ge0\), every \(a\in H^3_\sigma(\R^3)\), and every
\[
 0<\delta\le\delta_\nu(a)
 :=\min\{1,c_\nu(1+\norm a_{H^3})^{-2}\},
\]
the class \(\mathscr M_\nu^3(s,s+\delta;a)\) contains exactly one pair
\((w,q)\).  Its velocity satisfies
\[
 w\in C([s,s+\delta];H^3_\sigma)
 \cap L^2(s,s+\delta;H^4),
 \qquad
 w_t\in L^2(s,s+\delta;H^2).
\]
If \(a\in H^\infty_\sigma\), then
\[
 \partial_t^nw,\ \partial_t^nq
 \in C([s,s+\delta];H^m)
 \qquad(n,m\in\Nzero).
\]
For every finite \(b>s\), the class
\(\mathscr M_\nu^3(s,b;a)\) contains at most one pair.  Its compatible members
on strict subintervals of
\([s,s+\tau_{\nu,3}^{\max}(s,a))\) define the maximal shifted
pressure-normalized \(H^3\) mild history from \(a\).
\end{theorem}

\begin{proof}
By time translation it is enough to prove the assertions with \(s=0\).
Let \(0<\delta\le1\), let
\[
 z_t-\nu\Delta z=F,
 \qquad z(0)=a,
 \qquad
 a\in H^3_\sigma,\quad
 F\in L^2(0,\delta;H^2_\sigma),
\]
and set
\[
 A_0:=\norm a_{H^3},
 \qquad
 G:=\norm F_{L^2(0,\delta;H^2)}.
\]
The \(H^3\) energy identity is
\[
 \frac12\frac{\dd}{\dd t}\norm z_{H^3}^2
 +\nu\norm{\nabla z}_{H^3}^2
 =\operatorname{Re}\ip{A^2F}{A^4z}_{L^2}.
\]
Since
\(\norm z_{H^4}^2=\norm z_{H^3}^2+\norm{\nabla z}_{H^3}^2\),
Young's inequality gives
\[
 \frac{\dd}{\dd t}\norm z_{H^3}^2
 +\nu\norm{\nabla z}_{H^3}^2
 \le \nu\norm z_{H^3}^2+\nu^{-1}\norm F_{H^2}^2.
\]
Consequently,
\begin{equation}
 \norm z_{C([0,\delta];H^3)}
 \le e^{\nu/2}\bigl(A_0+\nu^{-1/2}G\bigr).
 \label{eq:restart-linear-CH3}
\end{equation}
Pairing \(A^2z_t-\nu\Delta A^2z=A^2F\) with
\(-\Delta A^2z\) gives
\[
 \frac12\frac{\dd}{\dd t}\norm{\nabla A^2z}_2^2
 +\nu\norm{\Delta A^2z}_2^2
 =\operatorname{Re}\ip{A^2F}{-\Delta A^2z}_{L^2}.
\]
Integration and Young's inequality yield
\begin{equation}
 \norm{\Delta z}_{L^2(0,\delta;H^2)}
 \le \nu^{-1/2}A_0+\nu^{-1}G.
 \label{eq:restart-linear-Laplacian}
\end{equation}
The pointwise multiplier inequality
\[
 \langle\xi\rangle^8
 \le2\bigl(\langle\xi\rangle^6
       +\abs\xi^4\langle\xi\rangle^4\bigr)
\]
and \(\delta\le1\) imply
\begin{equation}
 \norm z_{L^2(0,\delta;H^4)}
 \le\sqrt2\left(
 \norm z_{C([0,\delta];H^3)}
 +\norm{\Delta z}_{L^2(0,\delta;H^2)}
 \right).
 \label{eq:restart-linear-H4}
\end{equation}
The equation and \eqref{eq:restart-linear-Laplacian} give
\begin{equation}
 \norm{z_t}_{L^2(0,\delta;H^2)}
 \le\nu^{1/2}A_0+2G.
 \label{eq:restart-linear-time}
\end{equation}
Equations \eqref{eq:restart-linear-CH3}--\eqref{eq:restart-linear-time}
give
\begin{align*}
 &\norm z_{C H^3}+\norm z_{L^2H^4}+\norm{z_t}_{L^2H^2}\\
 &\quad\le
 e^{\nu/2}(A_0+\nu^{-1/2}G)
 +\sqrt2\left(
 e^{\nu/2}(A_0+\nu^{-1/2}G)
 +\nu^{-1/2}A_0+\nu^{-1}G\right)
 +\nu^{1/2}A_0+2G\\
 &\quad\le L_\nu(A_0+G).
\end{align*}
The last inequality follows from
\[
 \nu^{1/2}\le\frac{1+\nu}{2},
 \qquad
 \nu^{-1/2}\le\frac{1+\nu^{-1}}2,
 \qquad
 e^{\nu/2}\le e^{(\nu+\nu^{-1})/2}.
\]
Fourier truncation and passage to the limit justify the preceding linear
estimates for the stated data class.

For \(v,z\in H^3_\sigma\), Lemmas~\ref{lem:sobolev-product} and
\ref{lem:riesz-leray} give
\begin{align}
 \norm{\Pp\nabla\cdot(v\otimes z)}_{H^2}
 &\le C_{\rm q}\norm v_{H^3}\norm z_{H^3},
 \label{eq:restart-quadratic-bilinear}\\
 \norm{\Pp\nabla\cdot(v\otimes v-z\otimes z)}_{H^2}
 &\le C_{\rm q}(\norm v_{H^3}+\norm z_{H^3})
 \norm{v-z}_{H^3}.
 \label{eq:restart-quadratic-difference}
\end{align}
Indeed,
\[
 \norm{\nabla\cdot(v\otimes z)}_{H^2}^2
 \le3\sum_{i,j=1}^3\norm{v_i z_j}_{H^3}^2
 \le3(C_3^{\rm alg})^2\norm v_{H^3}^2\norm z_{H^3}^2.
\]

Define
\[
 X_\delta:=
 \left\{v\in C([0,\delta];H^3_\sigma)\cap L^2(0,\delta;H^4_\sigma):
 v_t\in L^2(0,\delta;H^2_\sigma)\right\},
\]
\[
 \norm v_{X_\delta}:=\norm v_{C H^3}
 +\norm v_{L^2H^4}+\norm{v_t}_{L^2H^2},
 \qquad
 X_\delta(a):=\{v\in X_\delta:v(0)=a\}.
\]
\(X_\delta\) is a Banach space with this norm, and \(X_\delta(a)\) is closed.
For \(A_0>0\), set \(R:=2L_\nu A_0\) and
\[
 \mathbb B_R:=\{v\in X_\delta(a):\norm v_{X_\delta}\le R\}.
\]
The set \(\mathbb B_R\) is nonempty: the zero-forcing heat trajectory
\(v^{\rm lin}(t)=e^{\nu t\Delta}a\) belongs to \(X_\delta(a)\), and the
linear estimate with \(F=0\) gives
\(\norm{v^{\rm lin}}_{X_\delta}\le L_\nu A_0<R\).
For \(v\in\mathbb B_R\), define
\[
 \Phi(v)(t):=e^{\nu t\Delta}a
 -\int_0^te^{\nu(t-s)\Delta}
 \Pp\nabla\cdot(v\otimes v)(s)\dd s.
\]
If \(\delta\le\delta_\nu(a)\), then
\[
 \delta^{1/2}\le
 \frac{1}{8C_{\rm q}L_\nu^2(1+A_0)}.
\]
\begin{samepage}
The linear estimate and \eqref{eq:restart-quadratic-bilinear} imply
\begin{align*}
 \norm{\Phi(v)}_{X_\delta}
 &\le L_\nu\bigl(A_0+C_{\rm q}\delta^{1/2}R^2\bigr)\\
 &\le L_\nu A_0+\frac{L_\nu A_0^2}{2(1+A_0)}
 <R.
\end{align*}
\end{samepage}
\begin{samepage}
For \(v,z\in\mathbb B_R\),
\begin{align*}
 \norm{\Phi(v)-\Phi(z)}_{X_\delta}
 &\le L_\nu C_{\rm q}\delta^{1/2}
 (\norm v_{C H^3}+\norm z_{C H^3})
 \norm{v-z}_{C H^3}\\
 &\le\frac{A_0}{2(1+A_0)}\norm{v-z}_{X_\delta}
 \le\frac12\norm{v-z}_{X_\delta}.
\end{align*}
\end{samepage}
Thus \(\Phi\) has a unique fixed point in \(\mathbb B_R\).  If \(A_0=0\),
the zero function is the fixed point.  Let \(b>0\) and let
\[
 (v,q_v),(z,q_z)\in\mathscr M_\nu^3(0,b;a).
\]
Set
\[
 M_b:=\norm v_{C([0,b];H^3)}+\norm z_{C([0,b];H^3)}.
\]
On every subinterval \(I\subset[0,b]\) of length \(h\le1\) on which the two
velocities have a common left trace, their shifted Duhamel identities and the
zero-data linear estimate give
\[
 \norm{v-z}_{C(I;H^3)}
 \le L_\nu C_{\rm q}h^{1/2}M_b
 \norm{v-z}_{C(I;H^3)}.
\]
A finite partition with
\(L_\nu C_{\rm q}h^{1/2}M_b<1\) proves \(v=z\) on \([0,b]\).
The pressure formula gives \(q_v=q_z\).  This proves the at-most-one
assertion for the full shifted mild class.  For the fixed point,
\eqref{eq:restart-quadratic-bilinear} places the forcing in
\(L^2(0,\delta;H^2)\); the linear estimates give
\[
 w\in L^2(0,\delta;H^4),
 \qquad
 w_t\in L^2(0,\delta;H^2).
\]
Lemma~\ref{lem:riesz-leray} gives the normalized pressure
\(q=R_iR_j(w_iw_j)\).

Fix \(m\ge3\) and assume \(a\in H^m_\sigma\).  Let \(S_k\) be Fourier
restriction to \(\{\abs\xi\le k\}\), set \(a_k=S_ka\), and replace the
fixed-point map above by
\[
 \Phi_k(v)(t):=e^{\nu t\Delta}a_k
 -\int_0^te^{\nu(t-s)\Delta}S_k\Pp\nabla\!\cdot
   (v\otimes v)(s)\dd s.
\]
The projection \(S_k\) is self-adjoint, contractive on every Sobolev space,
and commutes with the heat flow, \(\Pp\), and spatial derivatives.  These estimates that gave \(\Phi\) its contraction factor at most \(1/2\) apply to
\(\Phi_k\) on the closed radius-\(R\) ball in \(X_\delta(a_k)\).  Its fixed point
\(w^{(k)}=S_kw^{(k)}\) consequently exists on \([0,\delta]\) and satisfies
\begin{equation}
 \sup_k\norm{w^{(k)}}_{X_\delta}\le R.
 \label{eq:restart-cutoff-uniform-H3}
\end{equation}
Because its Fourier support is bounded, \(w^{(k)}\) belongs to every spatial
Sobolev space.  Self-adjointness of \(S_k\) removes the cutoff from the energy
pairing, so Lemma~\ref{lem:transport-commutator} gives
\[
 \frac12\frac{\dd}{\dd t}\norm{w^{(k)}}_{H^m}^2
 +\nu\norm{\nabla w^{(k)}}_{H^m}^2
 \le K_m\norm{\nabla w^{(k)}}_{L^\infty}
       \norm{w^{(k)}}_{H^m}^2.
\]
Set
\[
 C_{\nabla,3}:=(2\pi)^{-3/2}
 \left(\int_{\R^3}\abs\xi^2\langle\xi\rangle^{-6}\dd\xi\right)^{1/2}.
\]
Fourier inversion and Cauchy--Schwarz give
\(\norm{\nabla v}_\infty\le C_{\nabla,3}\norm v_{H^3}\).  From
\eqref{eq:restart-cutoff-uniform-H3},
\[
 \int_0^\delta\norm{\nabla w^{(k)}(s)}_\infty\dd s
 \le C_{\nabla,3}\delta R.
\]
Put
\[
 F_k(t):=\norm{w^{(k)}(t)}_{H^m}^2
 {}+2\nu\int_0^t\norm{\nabla w^{(k)}(s)}_{H^m}^2\dd s.
\]
The differential inequality above gives
\[
 F_k(t)\le\norm a_{H^m}^2
 {}+2K_m\int_0^t\norm{\nabla w^{(k)}(s)}_\infty F_k(s)\dd s.
\]
Gronwall's inequality now gives the uniform bound
\begin{equation}
 \sup_{0\le t\le\delta}\left(
 \norm{w^{(k)}(t)}_{H^m}^2
{}+2\nu\int_0^t\norm{\nabla w^{(k)}(s)}_{H^m}^2\dd s
 \right)
 \le e^{2K_mC_{\nabla,3}\delta R}\norm a_{H^m}^2.
 \label{eq:restart-cutoff-uniform-Hm}
\end{equation}
Since
\(
 \norm f_{H^{m+1}}^2=\norm f_{H^m}^2+\norm{\nabla f}_{H^m}^2
\), this also bounds \(w^{(k)}\) in \(L^2H^{m+1}\).  The cutoff equations
and Lemma~\ref{lem:sobolev-product} give, with a constant depending only on
\(m\),
\[
 \norm{\partial_tw^{(k)}}_{L^2H^{m-1}}
 \le \nu\norm{w^{(k)}}_{L^2H^{m+1}}
 +\sqrt3\,C_m^{\rm alg}\delta^{1/2}
   \norm{w^{(k)}}_{L^\infty H^m}^2,
\]
so the time derivatives are uniformly bounded in \(L^2H^{m-1}\).

It remains to identify the higher-order weak limit with the fixed point already
constructed in \(X_\delta(a)\).  The contraction estimate gives
\[
 \norm{w^{(k)}-w}_{X_\delta}
 \le2\norm{\Phi_k(w)-\Phi(w)}_{X_\delta}.
\]
The linear bound proved above controls the right-hand side by a constant times
\[
 \norm{(S_k-I)a}_{H^3}
 +\norm{(S_k-I)\Pp\nabla\!\cdot(w\otimes w)}_{L^2H^2},
\]
which tends to zero.  Hence \(w^{(k)}\to w\) in \(X_\delta\).  Weak
compactness in the three uniformly bounded higher-order spaces and uniqueness
of the distributional limit place
\[
 w\in L^\infty(0,\delta;H^m_\sigma)\cap L^2(0,\delta;H^{m+1}_\sigma),
 \qquad w_t\in L^2(0,\delta;H^{m-1}_\sigma).
\]
Lemma~\ref{lem:adjacent-trace} supplies the representative
\(w\in C([0,\delta];H^m_\sigma)\); its trace at zero is \(a\), since it is
the distributional trace of the already constructed \(C H^3\) solution.
Assume now that \(a\in H^\infty_\sigma\).  Set \(W_0=w\) and define
recursively
\[
 Q_n=R_iR_j\sum_{a=0}^n\binom naW_{a,i}W_{n-a,j},
\]
\[
 W_{n+1}=\nu\Delta W_n
 -\sum_{a=0}^n\binom na(W_a\cdot\nabla)W_{n-a}-\nabla Q_n.
\]
If \(W_0,\ldots,W_n\in C H^{m+2}\), then
Lemmas~\ref{lem:sobolev-product} and \ref{lem:riesz-leray} give
\[
 Q_n\in C H^{m+1},
 \qquad
 W_{n+1}\in C H^m.
\]
Since \(m\) is arbitrary, induction gives \(W_n,Q_n\in C H^m\) for every
\(n,m\in\Nzero\).  The equation first gives
\(W_0\in C^1H^m\) and \(\partial_tW_0=W_1\).  If
\(W_a=\partial_t^aw\) for \(0\le a\le n\), distributional differentiation
of \(q=R_iR_j(w_iw_j)\) and Lemma~\ref{lem:mixed-leibniz} give
\(\partial_t^nq=Q_n\).  Differentiation of the momentum equation then gives
\(\partial_tW_n=W_{n+1}\).  Induction gives
\[
 W_n\in C^1H^m,
 \qquad
 \partial_tW_n=W_{n+1},
 \qquad
 W_n=\partial_t^nw,
 \qquad
 Q_n=\partial_t^nq.
\]
Time translation proves the result for every \(s\ge0\).  Lemma
\ref{lem:shifted-mild-restriction} and the at-most-one assertion make all
nonempty classes below \(\tau_{\nu,3}^{\max}(s,a)\) compatible; their union
is the maximal shifted pressure-normalized \(H^3\) mild history.
\end{proof}
